
Reinforcement Learning from Human Feedback (RLHF) has become the dominant paradigm for aligning large language models with human preferences. Reward models trained to predict human judgments are evaluated on benchmarks such as RewardBench, which report aggregate accuracy metrics often exceeding 85\%. However, aggregate metrics do not reveal the structure of failures. When a reward model scores 85\% accuracy, what distinguishes the successful 85\% from the failed 15\%? More critically, within the successful majority, how many cases reflect genuine alignment versus confident misalignment---the model certain while humans disagree?

This paper introduces a decomposition framework that reveals patterns hidden by aggregate evaluation. Preference battles are classified along two orthogonal dimensions: human vote entropy (high when voters disagree, low when consensus) and reward model variance (high when uncertain, low when confident). Crossing these dimensions yields four alignment modes:

\begin{itemize}
\item \textbf{Mode 1 (Aligned-Confident)}: Low entropy, low variance---humans agree, reward model agrees
\item \textbf{Mode 2 (Aligned-Uncertain)}: Low entropy, high variance---humans agree, reward model uncertain
\item \textbf{Mode 3 (Misaligned-Confident)}: High entropy, low variance---humans disagree, reward model confident
\item \textbf{Mode 4 (Misaligned-Uncertain)}: High entropy, high variance---both uncertain
\end{itemize}

Mode 3 is the critical failure pattern. Reward models express high confidence on samples where human preferences diverge. If prevalent, this mode represents a systematic blind spot in RLHF training signal.

Analysis of 57,477 battles from the LMSYS Chatbot Arena finds Mode 3 constitutes 23.7\% of samples (95\% CI: 23.4--24.0\%, p < 0.001 against a 10\% threshold). This is not a fringe phenomenon but a substantial pattern affecting approximately one quarter of preference data.

Two mechanistic hypotheses were tested. First, the hypothesis that Mode 3 arises when responses are semantically divergent. The data falsifies this: Mode 3 response pairs show higher semantic similarity than Mode 1 pairs (Cohen's d = $-0.049)$, the opposite direction predicted. Second, the hypothesis that Mode 3 concentrates in subjective prompt types. The data shows negligible enrichment (ratio = 1.07 vs. predicted 1.5); Mode 3 pervades all task types.

\textbf{Contributions.} (1) The first empirical quantification of overconfident reward model misalignment in large-scale preference data, demonstrating that approximately 24\% of battles exhibit this failure pattern. (2) A $2\times 2$ mode decomposition (entropy $\times$ variance) as a diagnostic framework revealing patterns aggregate metrics miss. (3) Negative results ruling out semantic divergence and prompt subjectivity as explanatory mechanisms.
